\addchap{Comment lire ce manuel}

Ce manuel s'adresse à vous, ingénieur ou ingénieure, qui avez déjà livré des logiciels et
qui découvrez la démarche \emph{Lumière} et son plugin \textbf{Simetriik Forge}. Vous n'avez pas besoin
de tout lire dans l'ordre, ni de tout retenir : il est construit pour que vous pratiquiez vite, sur
un vrai projet, et que vous reveniez aux pages utiles quand le besoin se présente.

\section*{Ce que vous saurez faire}
\begin{itemize}
  \item cadrer un besoin en objectifs mesurables, et refuser une fonctionnalité qui n'en sert aucun ;
  \item bâtir un langage commun entre métier et code, et le faire respecter automatiquement ;
  \item écrire une spécification que le client peut valider et qu'un test peut vérifier ;
  \item consigner une décision d'architecture, chiffrée en FCFA, pour que personne ne la rediscute dans six mois ;
  \item déclencher la recette et la production sans improviser ;
  \item savoir \emph{pourquoi} un garde-fou vous refuse un commit, et quoi faire ensuite.
\end{itemize}

\section*{Un fil rouge : la plateforme Sika}
Tout au long du manuel, un même projet fictif sert d'exemple : \sika{}, une plateforme qui permet aux
groupes de tontine d'une association de commerçantes de cotiser par Mobile Money au lieu de
passer par un cahier et des espèces. Vous le suivrez de la première conversation avec la
présidente jusqu'à la revue post-production. Chaque concept est d'abord défini, puis
illustré sur \sika{}, puis proposé à votre propre projet.

\begin{situation}[Fiction]
\sika{}, l'association, les montants et les noms qui l'entourent sont inventés pour les besoins du
manuel. Les chiffres (coûts, taux, délais) illustrent un raisonnement ; ils ne sont pas des
recommandations et ne doivent pas être recopiés tels quels.
\end{situation}

\section*{Comment un chapitre est construit}
Les adultes apprennent mieux quand ils savent \emph{pourquoi} ils apprennent, partent de leur
expérience et appliquent tout de suite. Chaque chapitre d'étape suit donc le même rythme,
repérable à ses encadrés.

\begin{objectifs}[Objectifs]
Ce que vous serez capable de faire à la fin du chapitre, énoncé en verbes d'action.
\end{objectifs}
\begin{situation}[Sur le terrain : \sika{}]
Une scène du projet fil rouge, pour poser le problème avant la théorie.
\end{situation}
\begin{concept}[Concept]
La définition précise d'un mot de la démarche. Un seul sens, relu à chaque chapitre.
\end{concept}
\begin{piege}[Piège]
Une erreur fréquente, et comment l'éviter.
\end{piege}
\begin{vousjouez}[À vous de jouer]
Un exercice court à refaire sur \emph{votre} projet, tout de suite.
\end{vousjouez}
\begin{retenir}[À retenir]
L'essentiel en quelques lignes, à relire quand vous revenez au chapitre.
\end{retenir}

Chaque chapitre d'étape se termine par quelques questions d'auto-évaluation
(« Vérifiez vos acquis »), dont les réponses commentées figurent en annexe~\ref{app:corriges}.

\section*{Trois parcours de lecture}
\noindent\begin{tabularx}{\linewidth}{>{\raggedright\arraybackslash\bfseries\color{accent}}p{3.2cm}X}
\toprule
Je découvre & Lisez les chapitres \ref{ch:demarche} et \ref{ch:prise} dans l'ordre, puis les étapes une à une. Appliquez chaque « À vous de jouer » avant de passer au chapitre suivant. \\
Je connais Lumière & Parcourez le chapitre \ref{ch:prise}, puis allez à l'étape qui vous bloque. L'annexe~\ref{app:ref} liste tous les skills et scripts. \\
Je suis pressé(e) & Une petite fonctionnalité ? Chapitre \ref{ch:rapide}. Un commit refusé ? Annexe~\ref{app:depannage}. \\
\bottomrule
\end{tabularx}

\section*{Conventions}
\begin{itemize}
  \item Les commandes, chemins et noms de fichiers sont en \texttt{police à chasse fixe}. Les skills s'invoquent par leur nom, par exemple \skill{forge-cadrage}.
  \item Le code, les noms de fichiers et les routes sont en anglais ; les commentaires et la documentation, en français.
  \item Les sorties de scripts reproduites ici ont été obtenues en exécutant réellement les garde-fous sur un projet d'essai.
\end{itemize}

\begin{retenir}[Votre carnet de bord]
Gardez un projet réel sous la main et un carnet ouvert : notez, à chaque « À vous de jouer »,
ce qui s'applique à votre contexte et ce qui ne s'y applique pas. Une démarche ne devient la
vôtre que lorsque vous l'avez adaptée.
\end{retenir}
